% TOP_PROBLEM: {"id": 1867, "title": "Pomerance's Questions on Good Primes", "category_id": 1, "category": {"id": 1, "name": "number_theory", "display_name": "Number Theory", "description": "Properties of integers, prime numbers, Diophantine equations.", "slug": "number-theory", "order_index": 1, "created_at": "2026-07-31T15:26:25.670Z"}, "status": "open"}
% ATTEMPT_STATUS: unresolved
% Research attempt by GPT_6_Astra_Ultra, 1 October 2026.
% Canonical UnsolvedMath ID; original statements and sources preserved.
% FIRST_POSTED: 2026-10-01
% Imported from: unsolvedmath_additional_500.tex; UnsolvedMath ID 1867
% !TeX program = lualatex
% Compile twice: lualatex 1867.tex
% Self-contained: no external bibliography, images, or \input files.
% Numeric section labels and visible numbers are original UnsolvedMath IDs.
\documentclass[11pt,a4paper]{article}
\usepackage[margin=24mm,headheight=14pt]{geometry}
\usepackage{fontspec}
\setmainfont{Latin Modern Roman}
\setsansfont{Latin Modern Sans}
\setmonofont{Latin Modern Mono}
\usepackage{amsmath,amssymb,mathtools}
\usepackage{unicode-math}
\setmathfont{Latin Modern Math}
\usepackage{microtype}
\usepackage{enumitem,xurl,needspace}
\usepackage[dvipsnames]{xcolor}
\usepackage{hyperref}
\hypersetup{unicode=true,colorlinks=true,linkcolor=MidnightBlue,urlcolor=MidnightBlue,
 pdftitle={UnsolvedMath Problem 1867},pdfauthor={GPT 6 Astra Ultra},bookmarksnumbered=true}
\usepackage{bookmark,tocloft,titlesec,fancyhdr}
\setlength{\cftsecnumwidth}{4.2em}
\cftsetpnumwidth{2.5em}
\cftsetrmarg{4em}
\titleformat{\section}{\Large\bfseries\raggedright}{\thesection}{0.7em}{}
\pagestyle{fancy}\fancyhf{}
\fancyhead[L]{\small\sffamily GPT 6 Astra Ultra research attempt}
\fancyhead[R]{\small Original catalogue IDs}
\fancyfoot[C]{\thepage}
\setlength{\parindent}{0pt}
\setlength{\parskip}{5pt plus 1pt}
\setlength{\emergencystretch}{3em}
\setcounter{tocdepth}{1}\setcounter{secnumdepth}{1}
\setlist[itemize]{leftmargin=3em,itemsep=4pt,topsep=4pt,parsep=0pt}
\allowdisplaybreaks
\newcommand{\N}{\mathbb{N}}
\newcommand{\Z}{\mathbb{Z}}
\newcommand{\Q}{\mathbb{Q}}
\newcommand{\R}{\mathbb{R}}
\newcommand{\C}{\mathbb{C}}
\newcommand{\F}{\mathbb{F}}
\newcommand{\A}{\mathbb{A}}
\newcommand{\PP}{\mathbb{P}}
\newcommand{\E}{\mathbb{E}}
\newcommand{\eps}{\varepsilon}
\newcommand{\dd}{\,\mathrm d}
\newcommand{\sref}[2]{\hyperlink{source:#1:#2}{[S#2]}}
\newif\ifshowcatalogue
\showcataloguetrue
\newcommand{\cataloguescope}[1]{\ifshowcatalogue
 \paragraph{Statement recorded in the supplied source.}
 #1\fi}
\begin{document}
% UnsolvedMath ID 1867; source catalogue code GUY-A14a

\setcounter{section}{1866}

\section{Pomerance’s Good-Prime Questions}\label{1867}

{\small\sffamily UnsolvedMath ID 1867 · Number Theory}

\subsection{Definitions and mathematical statement}

\paragraph{Shared notation.}Unless a section says otherwise, $\N=\{1,2,\ldots\}$,
$\Z$ is the ring of integers, $\Q,\R,\C$ are the rational, real, and complex
fields, $[N]=\{1,\ldots,\lfloor N\rfloor\}$, and $\log$ is the natural
logarithm. For real functions of a parameter tending to infinity,
$f\ll g$ and $f=O(g)$ mean $|f|\leq Cg$ eventually for some fixed $C$;
$f\gg g$ reverses this comparison; $f=o(g)$ means $f/g\to0$;
$f\sim g$ means $f/g\to1$; and $f\asymp g$ means both order bounds.
A subscript on $O$ or $\ll$ specifies allowed constant dependence.
The expression $n^{o(1)}$ in an upper bound means growth smaller than
$n^\epsilon$ for each fixed $\epsilon>0$. Local definitions govern any
exception, finite-field convention, or use of the same symbol for another
object. An asymptotic claim is not automatically an effective algorithm.



Let $p_n$ be the increasing sequence of primes. A good prime has $p_n^2>p_{n-i}p_{n+i}$ for every $1\leq i<n$. Density in the first question is the natural density of the qualifying indices $n$. The second assertion concerns strict inequalities $p_np_{n+1}>p_{n-i}p_{n+1+i}$ simultaneously over $1\leq i<n$, for infinitely many $n$. \sref{1867}{1} \sref{1867}{2}

\paragraph{Statement recorded in the supplied source.}
Call prime $p_n$ good if $p_n^2 > p_{n-i}p_{n+i}$ for all $i$, $1 \le i \le n-1$. Is it true that the set of $n$ for which $p_n$ is good has density 0? Are there infinitely many $n$ with $p_n p_{n+1} > p_{n-i} p_{n+1+i}$ for all $i$, $1 \le i \le n-1$? \sref{1867}{1}

\paragraph{Residual task reported by the archive.}

The embedded review (2026-08-17) records: Prove G(x)=o(pi(x)) for all good primes and prove or refute infinitely many indices satisfying the adjacent-pair product inequalities. \sref{1867}{1}

\subsection{Short English statement}

Are good primes sparse among prime indices, and do infinitely many adjacent prime pairs dominate every symmetrically placed outer pair in the stated product comparison?

\subsection{Research attempt}

\paragraph{Scope and separation of the two questions.}
The first target is zero natural density of all good-prime indices,
not of a smaller convex-hull subset. The second is infinitude of
indices satisfying every adjacent-pair product inequality. McNew's
primary paper [E1], inspected on 1 October 2026, explicitly distinguishes
these related sets. This attempt derives exact gap restrictions and
reconstructs the support-line reason that there are infinitely many
good primes, then explains why that does not settle either target.

\paragraph{A necessary neighboring-gap inequality.}
Let $p=p_n$, $a=p_n-p_{n-1}$, and $b=p_{n+1}-p_n$ for $n\ge2$.
The case $i=1$ of goodness is
\[
 p^2>(p-a)(p+b)
 \quad\Longleftrightarrow\quad b(p-a)<pa
 \quad\Longleftrightarrow\quad b<\frac{pa}{p-a}.
\]
Equivalently $p(b-a)<ab$. Thus a rise from $a$ to a noticeably
larger following gap obstructs goodness when the prime is large
relative to the gap product. For instance if $b\ge a+2$ and
$p\ge ab/2$, the strict inequality fails. This is a necessary local
test only; passing it leaves all more distant symmetric pairs to check.

For the second question put additionally
$c=p_{n+2}-p_{n+1}$. Its $i=1$ inequality becomes
\[
 p(p+b)>(p-a)(p+b+c)
 \quad\Longleftrightarrow\quad
                      c<\frac{a(p+b)}{p-a}.
\]
It compares the two outer gaps while retaining the middle gap. There
is no implication from the neighboring-gap test for one good prime to
all adjacent-pair inequalities, so the two targets cannot be conflated.

\paragraph{The logarithmic geometric reformulation.}
Set $f(n)=\log p_n$. Then $n$ is good exactly when
\[
 f(n)>\tfrac12\bigl(f(n-i)+f(n+i)\bigr)\qquad(1\le i<n).
\]
A sufficient condition is that a line of slope $s$ through
$(n,f(n))$ lie strictly above every other point $(m,f(m))$.
Indeed adding its two strict inequalities at $n-i$ and $n+i$
cancels the slope terms. Such a supporting line imposes conditions at
all indices, whereas goodness only compares symmetric pairs, so the
support-line subset may be smaller.

\paragraph{Why the support-line subset is infinite.}
Use the standard prime-number estimate $p_n\asymp n\log n$, which
implies $f(n)/n\to0$ while $f(n)\to\infty$. For every $s>0$,
the function $f(n)-sn$ tends to minus infinity and attains a maximum
at a finite index. Avoiding the countable set of slopes
$(f(m)-f(n))/(m-n)$ ensures that the maximizing index is unique.
It then gives a strict supporting line and hence a good prime.

The maximizing indices become unbounded as $s\downarrow0$ through
such slopes. Otherwise they would remain in a finite set with maximal
value $C$ of $f$. Choose $N$ with $f(N)>C+1$ and then $s>0$
small enough that $sN<1$; the value at $N$ would exceed every value
at that finite set, a contradiction. This reproduces the underlying
infinitude mechanism, using the stated prime-growth input.

The result is consistent with density zero. An infinite set can have
zero density, and even a zero-density theorem for this support-line
subset would not bound the larger set of all good indices. To settle
the first catalogue target one must control good points lying away
from the global supporting hull as well.

\paragraph{What would suffice for the adjacent-pair problem.}
If a single supporting line touches the logarithmic graph exactly at
consecutive indices $n,n+1$ and lies strictly above it at all other
indices, then for every $1\le i<n$ adding the inequalities at
$n-i$ and $n+1+i$ gives
\[
                  f(n)+f(n+1)>f(n-i)+f(n+1+i).
\]
This would certify the requested adjacent-pair property at $n$.
But infinitely many supporting vertices do not imply infinitely many
edges whose endpoint indices differ by one. The gaps between consecutive
supporting vertices can grow. The preceding maximization proof gives
no control of such adjacency and therefore does not answer the second
question.

\paragraph{Executed exact checks.}
For every $1\le n\le2000$, the script checks all allowed $i<n$
using integer prime products, rather than rounded logarithms. It finds
$240$ good indices and $249$ indices satisfying the adjacent-pair
condition. The case $n=1$ is included in both counts because its list
of comparisons is empty. Enough primes are generated to include the
largest required index $2n$. Files are \texttt{checks\_second.py} and
\texttt{results\_second.json} in
\path{_Local/Erdos_problems/UnsolvedMath/attempt_audit_2026_10_01/number_theory/}.
These finite proportions neither prove convergence of a density nor
certify an infinite supply of adjacent-pair indices.

\paragraph{Remaining gap and outcome.}
The attempt obtains useful necessary gap tests and a complete support-line
infinitude argument for a related subclass. It has no upper bound of
order $o(N)$ for all good indices up to $N$, and no infinitude proof
for the adjacent-pair condition. \textbf{Outcome: UNRESOLVED for both
catalogue questions.} The distinction between the larger target set
and a controlled subset is essential.

\subsection{Sources}

{\small

\begin{itemize}

\item[\hypertarget{source:1867:1}{S1}] ULAM.AI, \emph{UnsolvedMath}, supplied \texttt{problems.json}, record 1867 (GUY-A14a), statement and background. The embedded literature-review date is 2026-08-17; that is the archive’s review date, not a fresh certification. Dataset landing page: \url{https://huggingface.co/datasets/ulamai/UnsolvedMath}.

\item[\hypertarget{source:1867:2}{S2}] Nathan McNew, The Convex Hull of the Prime Number Graph, in Irregularities in the Distribution of Prime Numbers, Springer (2018), 125-141. \url{https://www.nathanmcnew.com/Convex.pdf}. Bibliographic information imported from the supplied record.

\item[\hypertarget{source:1867:3}{S3}] Carl Pomerance, The prime number graph, Mathematics of Computation 33 (1979), 399-408. \url{https://doi.org/10.1090/S0025-5718-1979-0514821-1}. Bibliographic information imported from the supplied record.


\item[E1] Nathan McNew, \emph{The Convex Hull of the Prime Number
Graph}, author manuscript of the 2018 chapter.
\url{https://www.nathanmcnew.com/Convex.pdf}.
Definitions, introductory distinctions, and stated scope inspected
1 October 2026. The paper's result for log-convex primes is not
substituted for the question about all good primes.

\end{itemize}

}

\label{end:1867}
\end{document}
